% page 511 du fac-simile (image p-510 du scan)
% bloc 167.6 x 237.7 mm ; marge gauche 34.4 mm
\pgc{12.700mm}{0.000mm}{
\l{\cel{56}503.}
\l{}
\l{\sou{rosmaro}. (\sou{zool.}) Mamifero amfibia, karnavida, analoga a maro-}
\l{\cel{3}hundo. L.\cel{1}\sou{trichochus rosmarus}. - L.}
\l{}
\l{\sou{rosmarino}. (\sou{bot.}) Planto aromata, ek la familio \textquotedbl{}labiacei\textquotedbl{}, qua}
\l{\cel{3}kreskas precipue en maro. - L.\cel{1}\sou{rosmarinus}. - DEFIR.}
\l{}
\l{\sou{rospo}. (\sou{zool.}) Reptero batraka, kun gambi plu kurta kam olti di}
\l{\cel{3}rano, korpo grosa e kurta, sur qua abundegas la veruki, de qui}
\l{\cel{3}exudas odoro fetida. - L.\cel{1}\sou{bufo}. - I.}
\l{}
\l{\sou{rostar}. (\sou{trans}.)Koquar per fairo intensa (karno, pultruni) sur}
\l{\cel{3}spiso, en forno, sur grililo. - DEFI.}
\l{}
\l{\sou{rostro}. I. (\sou{che la antiqui}). Beko, quaze sporno, ye qua (kom armo)}
\l{\cel{3}onu provizis la pruo di la navi, por plufaciligar la asalt-}
\l{\cel{3}abordo di la navi enemika. - II. (\sou{zool.}) Longa prolonguro}
\l{\cel{3}di la nazo, che la elefanto, quan la animalo uzas quale ma-}
\l{\cel{3}nuon. (analoge pri la insekti).- III. (arkitekt.)\cel{2}Ornemento}
\l{\cel{3}di arkitekturo o di skulturo, forme di beko, di sporno.}
\l{}
\l{\sou{roto}. (\sou{tekn.}) Peco rigida, cirkla, qua rotacas cirkum axo}
\l{\cel{3}qua trairas lua centro, e pleas la rolo di motoro, generale}
\l{\cel{3}kun felgo e radii. - FIS.}
\l{}
\l{\sou{rotacar}. (\sou{netrans.) (aludante korpo)}. Movar su cirkum sua axo}
\l{\cel{3}fixa. - DEFIRS.}
\l{}
\l{\sou{rotano}. (\sou{bot.})\cel{2}Genero de palmiero, kun stipo longa e tenua.}
\l{\cel{3}- L.\cel{1}\sou{calamus rotang}. DEFIRS.}
\l{}
\l{\sou{rotifero}. (\sou{zool.)}\cel{2}Genero de animali mikroskopala, qui divenas}
\l{\cel{3}sen-mova kande vetero esas sika, e rimoveskas lor humidesko}
\l{\cel{3}di atmosfero. - EFIS.}
\l{}
\l{\sou{rotolar}. (\sou{netrans}.)Produktar bruiso (quale per frapi qui eventas}
\l{\cel{3}tre kurta-intertempe) kontinua, analoga ad olta de veturo qua}
\l{\cel{3}rulas rapide sur voyo ne-glata. -}
\l{}
\l{\sou{rotondo}. I. (\sou{edifico)}. Edifico cirklatra, quaze kronizata ye}
\l{\cel{3}kupolo. - II. (\sou{vest}o) Longa manteleto sen maniki qua envol-}
\l{\cel{3}vas la korpo. - DEFIRS.}
\l{}
\l{\sou{rotoro}. (\sou{tekn.}) Nomo di la parto nefixa, en la elektro-motori}
\l{\cel{3}kun fluo alternanta - opoze statoro, qua esas la parto fixa.-}
\l{}
\l{\sou{rotulo}. I\cel{1}\sou{(anat.).}\cel{2}Mikra osto plata, kun anguli kelke rondigita,}
\l{\cel{3}ye qua konsistas la parto avana di la genuo. - II. (\sou{tekn.})}
\l{\cel{3}Peco metala, sfero-forma, uzata kom artiko en la mashin-}
\l{\cel{2}organi qui bezonas povar movar su omna-direcione. - FIS.}
\l{}
\l{\sou{rovo}. (\sou{bot.}) Arbusto dornoza e reptera, ek la familio \textquotedbl{}rozacei\textquotedbl{},}
\l{\cel{3}qua kreskas en la hegi, la boski, e donas frukto (rovbero).}
\l{\cel{3}- L. rubus caesius,\cel{1}\sou{rubus fructicosus.}\cel{1}- I.}
}
